% !TEX root = ../main.tex
\section{Vertikal integrasjon og outsourcing}\label{thr:vertical-integration-outsourcing}

\paragraph{Definisjon.} \emph{Vertikal integrasjon} innebærer at et firma tar eierskap over aktiviteter i tilstøtende ledd i verdikjeden (oppstrøms leverandører eller nedstrøms distribusjon). \emph{Outsourcing} (market transactions) innebærer at disse aktivitetene anskaffes via kontrakter med uavhengige markedsaktører. Valget er en kjernebeslutning i \thref{boundaries-of-the-firm}{Boundaries of the Firm}. [SLIDE]

\subsection*{Forklaring}

BSZ kap.\ 8/19 diskuterer make-or-buy-beslutningen langs tre dimensjoner: transaksjonskostnader, insentiver og markedsmakt. [SLIDE]

\textbf{Fordeler ved outsourcing (markedstransaksjoner):}
\begin{enumerate}\setlength{\itemsep}{2pt}
\item \textbf{Markedsdisiplin og prismekanisme.} Konkurranseeksponering tvinger leverandøren til kostnadseffektivitet -- interne enheter mangler dette presset. [SLIDE]
\item \textbf{Spesialisering og skalafordeler.} Ekstern leverandør betjener mange kunder og oppnår stordriftsfordeler firmaet selv ikke kan matche. [SLIDE]
\item \textbf{Fleksibilitet.} Lettere å bytte leverandør enn å omstrukturere intern produksjon ved endret etterspørsel. [BOK]
\item \textbf{Sterkere insentiver.} Leverandøren er \emph{residual claimant} på sin profitt, mens en intern enhet typisk kompenseres med svakere insentiver (cost center). Jf.\ \thref{residual-claimants}{residual claimants}. [SLIDE]
\item \textbf{Unngå byråkratikostnader.} Hierarkisk styring bringer koordinasjon, monitoring og politiske kostnader som markedskontrakter slipper. [BOK]
\end{enumerate}

\textbf{Ulemper ved outsourcing (fordeler ved vertikal integrasjon):}
\begin{enumerate}\setlength{\itemsep}{2pt}
\item \textbf{Transaksjonskostnader.} Søke-, forhandlings-, overvåknings- og håndhevingskostnader kan overstige interne styringskostnader -- særlig ved \emph{spesifikke aktiva}. Jf.\ \thref{costly-contracting}{costly contracting}. [SLIDE]
\item \textbf{Hold-up-problemet.} Når relasjonen krever spesifikke investeringer, risikerer parten med svakest forhandlingsposisjon at motparten reforhandler ex post. Integrasjon eliminerer hold-up. Se \thref{hold-up-problem}{hold-up-problemet}. [SLIDE]
\item \textbf{Free-rider-problemet.} Uavhengige distributører kan underinvestere i service/merkevare fordi gevinsten deles med andre. Integrasjon internaliserer eksternaliteten. Se \thref{free-rider-problem}{free-rider-problemet}. [BOK]
\item \textbf{Double markup.} Når to suksessive monopolister priser uavhengig, blir sluttprisen for høy og kvantum for lavt -- vertikal integrasjon fjerner den ene markupen. Se \thref{double-markup-problem}{double markup-problemet}. [SLIDE]
\item \textbf{Koordinasjon og informasjonsflyt.} Intern koordinasjon kan være enklere enn å formidle \thref{specific-vs-general-knowledge}{spesifikk viten} via prismekanismen. [BOK]
\item \textbf{Beskyttelse av immaterielle rettigheter.} Outsourcing krever deling av teknologi/prosesser -- integrasjon beskytter trade secrets. [BOK]
\item \textbf{Market power.} VI kan brukes strategisk for å kontrollere tilgang til inputs eller distribusjon -- se \thref{vertical-integration-market-power}{VI og market power}. [SLIDE]
\end{enumerate}

\subsection*{Kjerneforutsetninger}
\begin{itemize}
\item Make-or-buy er et dikotomt valg (i praksis finnes mellomformer: franchising, JV, allianse).
\item Transaksjonskostnader og byråkratikostnader er sammenlignbare.
\item Kontraktsufullstendighet er grunnen til at markedet feiler -- med \thref{costless-contracting}{costless contracting} ville outsourcing alltid dominert.
\item Firmaets interne organisasjon er \emph{ikke} gratis -- agentkostnader internt erstatter transaksjonskostnader i markedet.
\end{itemize}

\subsection*{Muligheter}
\begin{itemize}
\item Forklarer verdikjedebeslutninger i de fleste industrier (bilindustri, energi, teknologi, shipping).
\item Predikerer outsourcing-trenden: når teknologi senker søke-/transaksjonskostnader (digitale plattformer, standardisering), øker outsourcing.
\item Predikerer integrasjon ved: høy asset specificity, dårlig leverandørmarked, store eksternaliteter.
\end{itemize}

\subsection*{Begrensninger og kritikk}
\begin{itemize}
\item Mellomformer (franchising, joint ventures, strategiske allianser) passer ikke ren dikotomi.
\item Ignorerer dynamikk: integrasjon i dag kan bli en byrde i morgen (lock-in i vertikal struktur).
\item Kulturelle og ledelsesmessige kapasitetshensyn: integrasjon krever lederkompetanse i nye verdikjedeledd.
\item Strategiske/politiske motiver (nasjonalt eierskap, regulering) overstyrer transaksjonskostnadsvurdering.
\end{itemize}

\subsection*{Modell / illustrasjon}

\begin{figure}[H]
  \centering
  \includegraphics[width=0.5\linewidth]{BSZ_kap_8_2026/by_slide/slide_009/BSZ_kap_8_2026_slide_009_asset_02_image5.jpeg}
  \caption{Spesialisert produksjonsutstyr (robotarm-fabrikk): make-or-buy-beslutningen avhenger av aktivaspesifisitet -- jo mer relasjonsspesifik investering, desto sterkere argument for vertikal integrasjon for å redusere hold-up-risiko og transaksjonskostnader. (BSZ kap.\ 8, slide 9) [SLIDE]}
\end{figure}

\begin{center}
\renewcommand{\arraystretch}{1.15}
\begin{tabular}{@{}p{4.0cm}p{5.0cm}p{5.0cm}@{}}
\toprule
 & \textbf{Vertikal integrasjon (make)} & \textbf{Outsourcing (buy)} \\
\midrule
Transaksjonskostnader & Lave (intern styring) & Kan være høye (søk, forhandling, håndhevelse) \\
Byråkratikostnader & Kan være høye (hierarki) & Lave (markedsdisiplin) \\
Insentiver & Svakere (\thref{agency-costs}{agency costs}) & Sterkere (residual claimant) \\
Fleksibilitet & Lav (lock-in i kapasitet) & Høy (bytte leverandør) \\
Koordinasjon & Direkte styring & Via kontrakt/pris \\
Hold-up-risiko & Eliminert & Kritisk ved spesifikke aktiva \\
Market power & Mulig (kontroll verdikjede) & Begrenset \\
\bottomrule
\end{tabular}
\end{center}
\textit{Tabell: Sammenligning av vertikal integrasjon og outsourcing langs sentrale dimensjoner. [SLIDE/BOK]}

\subsection*{Eksempler}
\begin{itemize}\setlength{\itemsep}{2pt}
\item \textbf{Apple (outsourcing + kontroll):} Apple designer selv (integrert) men produserer hos Foxconn/TSMC (outsourcet). Fordi chip-design er svært spesifikt men produksjonsteknologi er standardisert, har Apple valgt en hybrid: eierskap til IP, markedskontrakt for produksjon. [EKS]
\item \textbf{Equinor og brønnservice:} Equinor outsourcer brønntjenester (Schlumberger, Halliburton) fordi kompetansen er generell og leverandørmarkedet er konkurrerende. Men feltspesifikt utstyr kan skape hold-up -- løst med flerårige rammeavtaler med opsjonsklausuler. [EKS]
\item \textbf{Mærsk og APM Terminals:} Mærsk integrerte nedstrøms i terminal-operasjoner for å sikre havnekapasitet -- et \emph{bottleneck}-aktivum med høy spesifisitet. [CASE]
\end{itemize}

\subsection*{Kobling til søyle og andre teorier}
Søyle: \SOne\ (beslutningsrett over verdikjedeledd), \SThree\ (insentivdesign for interne vs.\ eksterne enheter). Relaterte: \thref{boundaries-of-the-firm}{Boundaries of the Firm}, \thref{costly-contracting}{Costly contracting}, \thref{hold-up-problem}{Hold-up-problemet}, \thref{free-rider-problem}{Free-rider-problemet}, \thref{double-markup-problem}{Double markup-problemet}, \thref{vertical-integration-market-power}{VI og market power}, \thref{agency-costs}{Agency costs}, \thref{specific-vs-general-knowledge}{Spesifik vs.\ generell viden}, \thref{transfer-pricing}{Transfer pricing}, \thref{accounting-vs-economic-profit}{Accounting vs.\ economic profit}.

\paragraph{Brukt i spørsmål:} Q11, Q12.
